\begin{table}[t]
\caption{The routing-policy portfolio. Every policy is an established family;
none is introduced here. The last column states a hypothesis fixed before the
experiment.}
\label{tab:portfolio}
\centering\small
\begin{tabularx}{\textwidth}{@{}l l l X X@{}}
\toprule
& Objective & Information used & Intended strength & Hypothesised failure mode\\
\midrule
P1 & $\min_p L(p)$ & none &
  needs no infrastructure and cannot oscillate &
  loads the shortest corridor past its capacity\\[2pt]
P2 & $\min_p \hat\mu_p(t-\Delta)$ & lagged link travel times &
  follows congestion as it forms &
  moves the whole guided cohort together on stale information\\[2pt]
P3 & $\min_{p\in\mathcal{A}} \max_{e\in p}\rho_e$ &
  lagged travel times and flows, plus its own recent assignments &
  spreads load when capacity is unevenly distributed &
  diverts to longer corridors when no capacity shortage exists\\[2pt]
P4 & $\min_p [\hat\mu_p + \lambda\hat\sigma_p]$ & lagged link travel times &
  avoids corridors whose travel time is volatile &
  pays a detour for variance that carries no risk\\
\bottomrule
\end{tabularx}
\end{table}